\section{Part 1：深度 RL 的开放问题与前沿}

Chelsea Finn 在本讲中总结了深度 RL 领域的开放问题和前沿方向，并提供了深度 RL 研究的实践建议。

\subsection{问题定义层面的挑战}

\subsubsection{奖励设计}

\begin{importantbox}{奖励设计的根本困难}
在现实应用中，设计正确的奖励函数往往比解决 RL 问题本身更困难：
\begin{itemize}
    \item \textbf{奖励稀疏}：大多数时间步没有有意义的反馈
    \item \textbf{奖励误指定}：设计的奖励可能导致非预期行为（reward hacking）
    \item \textbf{多目标冲突}：多个奖励目标之间的权重很难调节
\end{itemize}
\end{importantbox}

\begin{warningbox}{Reward Hacking}
当奖励函数不完美时，RL 智能体会找到"合法但非预期"的方式来最大化奖励。经典案例：
\begin{itemize}
    \item 赛车游戏中收集奖励币而非完成赛道
    \item 机器人通过抖动而非行走来最大化速度奖励
    \item LLM 通过重复输出来最大化长度奖励
\end{itemize}
Reward hacking 是 RL 对齐（alignment）问题的核心挑战之一。
\end{warningbox}

\begin{figure}[H]
\centering
\includegraphics[width=0.9\textwidth]{slides-images/slide-05.jpg}
\caption{课程用聊天机器人、机器人折叠和推荐系统说明：一旦任务缺少可靠奖励或客观可验证指标，RL 的问题设定本身就会变成研究难点。}
\end{figure}

\subsubsection{状态与动作空间设计}

如何选择合适的状态表示和动作空间也是一个重要的设计问题：
\begin{itemize}
    \item 状态应该包含足够的信息（满足 Markov 性），但也不能太高维
    \item 动作空间的粒度影响学习效率
    \item 在 LLM 中，是 token 级还是 sentence 级的动作空间
\end{itemize}

\subsubsection{环境设计与基准}

\begin{knowledgebox}{好的 RL 基准的标准}
一个有用的 RL 基准应该：
\begin{itemize}
    \item 捕捉真实应用中的核心挑战
    \item 可以快速迭代（计算开销适度）
    \item 评估指标有意义且可靠
    \item 与真实世界性能有相关性
\end{itemize}
当前的挑战是：很多流行的 RL 基准与真实应用的需求脱节。
\end{knowledgebox}

\subsection{方法层面的挑战}

\subsubsection{样本效率}

尽管已有很多改进，深度 RL 的样本效率仍然远低于人类学习。

\subsubsection{泛化能力}

\begin{importantbox}{RL 中的泛化问题}
当前 RL 方法的泛化能力有限：
\begin{itemize}
    \item 在训练环境外的新环境中性能急剧下降
    \item 对奖励函数的微小变化过度敏感
    \item 难以将在一个任务上学到的知识迁移到相关任务
\end{itemize}
提升 RL 泛化能力是实现真实世界部署的关键。
\end{importantbox}

\subsubsection{多智能体与博弈}

多智能体 RL 面临额外的挑战：
\begin{itemize}
    \item 环境非平稳（其他智能体的策略在变化）
    \item 协调与竞争的平衡
    \item 通信协议的学习
\end{itemize}

\begin{figure}[H]
\centering
\includegraphics[width=0.9\textwidth]{slides-images/slide-12.jpg}
\caption{讲者把 video generation / world model 放进方法前沿中讨论：它们提供了丰富先验，但在策略分布外使用时又会因为物理误差和分布外动作而失真。}
\end{figure}

\subsection{部署与评估层面的挑战}

\subsubsection{安全性与可信度}

\begin{warningbox}{部署 RL 系统的风险}
将 RL 系统部署到真实世界时面临多重风险：
\begin{itemize}
    \item 分布外（out-of-distribution）的状态可能导致灾难性行为
    \item RL 策略的决策过程难以解释
    \item 对抗攻击可能利用策略的弱点
    \item 长期影响难以预测
\end{itemize}
\end{warningbox}

\subsubsection{可重复性}

深度 RL 实验的可重复性问题一直备受关注：
\begin{itemize}
    \item 对随机种子敏感
    \item 超参数调优开销大
    \item 不同实现之间的差异可能很大
\end{itemize}

\subsection{本章小结}

深度 RL 面临从问题定义到方法再到部署的全方位挑战。奖励设计、泛化能力和安全性是最核心的开放问题。

